\documentclass[11pt,a4paper]{article}
\emergencystretch=3em
\usepackage{fontspec}
\setmainfont{Times New Roman}
\usepackage[margin=2.5cm]{geometry}
\usepackage{setspace}
\onehalfspacing

\begin{document}

\begin{flushright}
August 14, 2026
\end{flushright}

\vspace{12pt}

\noindent Dear Editor,

\vspace{12pt}

We submit for consideration at \textit{The Lancet Global Health} our
manuscript (Article) entitled ``\textbf{Domestic Government Health Expenditure, Out-of-Pocket Spending, and
Healthy Life Expectancy in 190 Countries, 2000--2022.}''

\textbf{What we found.} In a panel of 190 countries through 2022 (complete-case analyses on 189), we
separate between-country development gradients from within-country changes
in domestic government health expenditure (GHE). the pooled within-country association was null (TWFE
$\beta = -0.12$, $p = 0.077$; one-sided exclusion testing rules out
effects above $+0.10$ HALE-years per percentage point of GDP); the
low-income subgroup showed a positive but statistically uncertain signal
(wild cluster bootstrap $p = 0.065$); and GHE was consistently associated
with lower out-of-pocket shares ($\approx$11\% relative reduction per
percentage point of GDP, $p < 0.001$).

\textbf{Relevance to \textit{The Lancet Global Health}.} This Article
addresses a central policy question for universal health coverage: whether
increases in domestic public health financing are associated with
measurable gains in healthy life expectancy and financial protection,
particularly in low-income settings. Our estimates suggest that the
GHE-to-GDP ratio alone is not a strong predictor of short-to-medium-term
HALE change overall, while public financing is consistently associated with
lower out-of-pocket expenditure shares. The positive but uncertain
low-income signal is, in our view, hypothesis-generating and directly
relevant to forthcoming domestic financing commitments. Recent UHC monitoring documents slowing progress and persistent service gaps in low-income countries alongside continuing
financial-protection challenges from out-of-pocket spending, while
near-term ODA projections indicate further cuts by major donors, making evidence on the returns to domestic public financing more urgent.

\textbf{Methods and transparency.} Analyses were implemented through
parallel pipelines with consistent conclusions, reporting follows STROBE
and GATHER, and event studies, synthetic control case studies, and pathway
decompositions are reported as exploratory. A complete anonymised
replication package is available to editors and reviewers on request and
will be deposited in Zenodo with a DOI. AI-assisted tools were used for
language editing and structural suggestions; all output was reviewed and
edited by the authors, who take full responsibility for the content, and
no AI tool was used for data generation or statistical analysis.

\noindent 190-country panel; 4,304 complete-case country-years (189
countries; 2000--2022); $\sim$3335 words (main text); 3 tables, 3 figures; 20
supplementary tables, 14 supplementary figures. This study used publicly available, de-identified aggregate
country-level data and did not require institutional review board approval.

\noindent This study was supported by the National Natural Science
Foundation of China (82403759 to CX; 82403616 to YW). CX and YW provided
the research funding through these grants and participated in the data
analysis. Competing interests: none declared.

\vspace{12pt}

\noindent We confirm that this manuscript has not been published elsewhere and
is not under consideration by another journal. All authors have approved the
manuscript and agree with its submission to \textit{The Lancet Global Health}.

\vspace{18pt}

\noindent Sincerely,

\vspace{12pt}

\noindent Tao Zhu, MD, PhD (corresponding author)\\
Department of Obstetrics and Gynecology\\
Tongji Hospital, Tongji Medical College\\
Huazhong University of Science and Technology\\
Wuhan, China\\
Tel: +86-27-8366-2663\\
Email: zhutao@tjh.tjmu.edu.cn

\end{document}
